\chapter{Progettazione e architettura}
\label{chap:progettazione-architettura}

% ==================================================
% 4.1 Visione architetturale
% ==================================================

\section{Visione architetturale del sistema}
\label{sec:visione-architetturale}

% Obiettivo:
% spiegare le scelte architetturali come risposta ai
% requisiti, non come elenco di cartelle.
%
% Contenuti attesi:
% - applicazione web e monorepo;
% - separazione tra app e pacchetti riutilizzabili;
% - motivazione della modularità;
% - relazione tra dominio, UI, persistenza, geolocalizzazione
%   e servizi esterni.
%
% Verificare sempre contro il codice corrente prima di
% descrivere pacchetti, feature o responsabilità.

% ==================================================
% 4.2 Stack tecnologico
% ==================================================

\section{Stack tecnologico e criteri di scelta}
\label{sec:stack-tecnologico}

% Trattare tecnologie ricorrenti usando le macro definite
% in config/commands.tex quando disponibili.
%
% Tecnologie da discutere solo se confermate dal repository:
% - \nextjs, \reactjs, \typescript e \nx;
% - \postgresql, \postgis e \prisma;
% - \authjs;
% - \maplibre;
% - provider opzionali per geocoding, storage e OCR.
%
% Collegare ogni tecnologia al problema risolto.

% ==================================================
% 4.3 Modello dei dati
% ==================================================

\section{Modello dei dati e persistenza}
\label{sec:modello-dati}

% Derivare questa sezione da packages/db/prisma/schema.prisma.
%
% Contenuti attesi:
% - utenti e verifica email;
% - libri e metadati bibliografici;
% - immagini e thumbnail;
% - coordinate private e pubbliche;
% - statistiche;
% - richieste di consultazione/prestito;
% - ruoli e vincoli principali.
%
% Non inventare entità non presenti nello schema corrente.

% ==================================================
% 4.4 Geolocalizzazione e privacy
% ==================================================

\section{Progettazione della geolocalizzazione e tutela della posizione}
\label{sec:progettazione-geolocalizzazione-privacy}

% Spiegare la scelta di separare coordinate precise e
% coordinate pubbliche approssimate.
%
% Contenuti attesi:
% - geocoding dell'indirizzo;
% - approssimazione delle coordinate;
% - uso delle coordinate pubbliche per mappe e discovery;
% - query spaziali e ordinamento per distanza;
% - trade-off tra utilità e protezione.

% ==================================================
% 4.5 Integrazione con servizi esterni
% ==================================================

\section{Servizi esterni e strategie di fallback}
\label{sec:servizi-esterni-fallback}

% Descrivere solo integrazioni realmente presenti:
% - geocoding/autocomplete;
% - Open Library;
% - OCR opzionale;
% - storage immagini;
% - invio email.
%
% Evidenziare dove il sistema resta funzionante in modalità
% locale o degradata e dove richiede configurazione esplicita.
